\documentclass[10pt]{article}

\usepackage{microtype,amsmath,amssymb,amsthm,booktabs}
\usepackage[margin=0.85in]{geometry}
\usepackage[colorlinks=true,linkcolor=blue,urlcolor=blue,citecolor=blue]{hyperref}
\usepackage{bookmark}
\newtheorem{theorem}{Theorem}
\newtheorem{lemma}[theorem]{Lemma}
\newtheorem{corollary}[theorem]{Corollary}
\newcommand{\dd}{\mathrm d}
\newcommand{\Tr}{\operatorname{Tr}}
\newcommand{\artanh}{\operatorname{artanh}}
\setlength{\emergencystretch}{2em}
\title{Spectral Tails and Exact Concentration Logits}
\author{Kirtus Justus (Theoretical Physicist)\\
\small Port Coquitlam, British Columbia, Canada}
\date{October 2026}
\hypersetup{pdftitle={Spectral Tails and Exact Concentration Logits},pdfauthor={Kirtus Justus (Theoretical Physicist)}}
\usepackage{fontspec,unicode-math,mathtools}
\setmainfont{DejaVu Serif}\setsansfont{DejaVu Sans}\setmonofont{DejaVu Sans Mono}\setmathfont{Latin Modern Math}
\begin{document}
\maketitle
\begin{abstract}
For the half-filled finite sine kernel, we express the concentration logit
as a convergent sum over the spectrum of a singular differential operator
after its first $N$ eigenvalues have been removed. The finite Fourier
subspace is proved to exhaust those first $N$ modes. A Green-kernel trace
identity then determines the logarithmic derivative of the deflated
connection coefficient, with no undetermined entire factor. The resulting
matrix identity is an absolutely convergent odd-power series with positive
coefficients. It gives exact sign and monotonicity statements and explicit
spectral-tail truncation bounds at fixed $N$. Formal center coefficients
are compared with this representation. Their large-$N$ interpretation is developed in the reciprocal-trace companions, and the fixed-interior quantum-period remainder is proved separately in \emph{Uniform Bulk Modular Asymptotics and the Leading Wilson Entropy Law}.
\end{abstract}

\section{Finite reference problem}
Fix an integer $N\ge1$. For $0\le j,k<N$, let
\begin{equation}
B_{jk}=\begin{cases}1/2,&j=k,\\
\sin[\pi(j-k)/2]/[\pi(j-k)],&j\ne k,\end{cases}
\qquad
H_{j-1,j}=H_{j,j-1}=\frac{j(N-j)}2,\quad H_{jj}=0.
\label{eq:matrices}
\end{equation}
This is the classical discrete concentration geometry used in the
free-fermion entanglement problem \cite{eisler}. We use
\begin{equation}
\mathcal P_N=\log B-\log(I-B),\qquad K_N=-\mathcal P_N,
\label{eq:modular}
\end{equation}
with natural logarithms. The kernel is strictly between zero and one:
its quadratic form and that of its complement are integrals of the
squared absolute value of a nonzero finite Fourier polynomial over
nonempty arcs. No physical clock or energy scale is assumed.

Put $a=\pi/2$, $I=(-a,a)$, $\alpha=(N^2-1)/4$, and
\begin{equation}
U_v(x)=\sum_{j=0}^{N-1}v_j e^{i[j-(N-1)/2]x}.
\end{equation}
Direct Fourier summation gives
\begin{equation}
\mathcal L_N U_v=U_{Hv},\qquad
\mathcal L_N u=(\cos x\,u')'+\alpha\cos x\,u.
\label{eq:Fourier}
\end{equation}
The outer Fourier coefficients vanish because
$\alpha-m(m\pm1)=0$ at the appropriate endpoints of the frequency set.
Integration by parts on $I$, where $\cos x$ vanishes at the endpoints,
proves $HB=BH$. Since $H$ has positive bonds, it has simple spectrum.
Its real eigenvectors therefore give common eigenpairs
$Hv=\beta v$, $Bv=\zeta v$, with $0<\zeta<1$.

The diagonal matrix $E_{jj}=(-1)^j$ obeys
\begin{equation}
EHE=-H,\qquad EBE=I-B.
\label{eq:duality}
\end{equation}
Thus the spectrum of $H$ is symmetric. Write $h_N=\|H\|$.
The case $N=1$ has $h_N=0$ and $\mathcal P_N=0$.

\section{The differential spectrum and its first finite block}
Let $A_N$ be the self-adjoint Friedrichs operator associated with the
closure on $C_c^\infty(I)$ of
\begin{equation}
\mathfrak a_N[f]=\int_I \cos x\,|f'|^2\,\dd x
-\alpha\int_I\cos x\,|f|^2\,\dd x.
\label{eq:form}
\end{equation}
The differential expression is $A_N=-\mathcal L_N$.
Its eigenfunctions use the regular, logarithm-free solution at both
endpoints. We justify this choice below rather than selecting a different
self-adjoint extension implicitly. The role of this condition in
singular Sturm--Liouville problems is classical \cite{kaper,littlejohn}.

\begin{lemma}[Compact spectrum and elementary bounds]
The spectrum of $A_N$ is simple and discrete. In increasing order,
with zero-based indices, write it as $\mu_0<\mu_1<\cdots$.
For $c=2/\pi^2$,
\begin{equation}
c\,k(k+1)-\alpha\ \le\ \mu_k\ \le\ \frac12 k(k+1).
\label{eq:bounds}
\end{equation}
In particular its resolvent is trace class away from the spectrum.
\end{lemma}
\begin{proof}
For $|x|<a$,
\[
\frac{\cos x}{a^2-x^2}
=\frac12\frac{\sin[(a-x)/2]}{(a-x)/2}
          \frac{\sin[(a+x)/2]}{(a+x)/2}
\]
lies between $2/\pi^2$ and $1/2$.
Thus the positive derivative form has the same form domain and equivalent
form norm as that of the rescaled Friedrichs Legendre operator
$-[(a^2-x^2)f']'$, whose eigenvalues are $k(k+1)$ \cite{littlejohn}.
That form domain embeds compactly into $L^2(I)$.
The potential $-\alpha\cos x$ is bounded between $-\alpha$ and zero.
The min--max principle now gives compactness and \eqref{eq:bounds}.

At either endpoint, the two local solution types of the eigen-equation
are analytic and analytic times a logarithm plus analytic.
The logarithmic type has infinite derivative energy
$\int_0^\epsilon r\,|1/r|^2\,\dd r$ and is excluded by the form domain.
The regular type belongs to the closure: a cutoff from zero to one
over $\epsilon^2<r<\epsilon$, linear in $\log r$, has derivative energy
$O(1/|\log\epsilon|)$ and vanishing $L^2$ error.
For these regular solutions the endpoint flux vanishes, so the
eigen-equation also holds in the operator domain.
The regular solution at the left endpoint is unique up to a constant;
hence each eigenvalue is simple.
This argument concerns eigenfunctions, not arbitrary form-domain
functions, which need not be bounded at the endpoints.
Finally \eqref{eq:bounds} implies summability of the resolvent eigenvalues.
\end{proof}

\begin{theorem}[Exact selection of the first $N$ modes]
The eigenvalues of $A_N$ on its first $N$ modes are exactly the
negatives of the eigenvalues of $H$. Consequently
\begin{equation}
\mu_{N-1}=h_N,\qquad \mu_N>h_N\ge0.
\label{eq:gap}
\end{equation}
\end{theorem}
\begin{proof}
Let $V_N$ be the real-valued part of the centered Fourier space
spanned by the functions in \eqref{eq:Fourier}. Its real dimension is
$N$; it is invariant under $A_N$, and its eigenfunctions are regular
at both endpoints. Thus it supplies $N$ distinct eigenvalues.

A real continuous function orthogonal to $V_N$ must have at least $N$
sign changes if it has only finitely many zeros.
Indeed, if its sign-change points are $z_1,\ldots,z_r$ with $r<N$, then
\[
q(x)=\prod_{\ell=1}^{r}\sin\frac{x-z_\ell}{2}\,
                 \cos^{\,N-1-r}\frac{x}{2}
\]
belongs to $V_N$. On $I$, it has exactly those sign changes.
The product of the function and $q$ has a fixed sign and nonzero integral,
contradicting orthogonality.

The eigenfunction of index $k$ has at most $k+1$ nodal domains.
For completeness, its restrictions to nodal domains lie in the closed
form domain and have disjoint supports, each with Rayleigh quotient
$\mu_k$. If there were more than $k+1$ such pieces, their span would
contain a nonzero vector orthogonal to the first $k+1$ eigenfunctions.
Its quotient would be both $\mu_k$ and at least $\mu_{k+1}$, a
contradiction. Interior zeros are simple, and endpoint regularity
prevents an accumulation there, so the sign-change argument applies.

An eigenfunction of index $k<N$ outside $V_N$ would be orthogonal to
the invariant finite space, and would have at least $N$ sign changes,
contradicting the nodal bound. Hence $V_N$ contains the first $N$ modes.
Symmetry \eqref{eq:duality} identifies the largest selected value as
$h_N$, and simplicity gives the strict inequality in \eqref{eq:gap}.
\end{proof}

\section{Connection coefficients and a resolvent trace}
Let $u(x;b)$ solve
\begin{equation}
(\cos x\,u')'+(\alpha\cos x-b)u=0,\qquad u(-a;b)=1,
\label{eq:u}
\end{equation}
with the regular left endpoint condition. Let $v(x;b)=u(-x;b)$.
At the right endpoint define $C_N(b)$ by
\[
u(x;b)=C_N(b)v(x;b)\log(a-x)+\text{an analytic function}.
\]
The regular right factor has value one at $a$. The normalized Frobenius
recursion and continuation show that $C_N$ is entire in $b$, and
its constant Wronskian is
\begin{equation}
W=\cos x\,(uv'-u'v)=C_N(b).
\label{eq:W}
\end{equation}
In particular $C_N(b)=-2u(0;b)u'(0;b)$.
Its zeros are precisely $b=-\mu_k$. They are simple:
at a real zero, reflection gives $u(a)=h=\pm1$, and parameter
differentiation yields
\begin{equation}
\int_I u^2\,\dd x=-h\,C_N'(b)>0.
\label{eq:norm}
\end{equation}
There are no nonreal zeros, because a regular solution at both endpoints
would give an eigenvalue of the self-adjoint operator.

\begin{lemma}[Trace identity without a normalization ambiguity]
For $b$ outside the zeros of $C_N$,
\begin{equation}
\frac{C_N'(b)}{C_N(b)}
=\Tr(A_N+b)^{-1}
=\sum_{k=0}^\infty\frac1{\mu_k+b}.
\label{eq:trace}
\end{equation}
\end{lemma}
\begin{proof}
The Green kernel of $A_N+b$ is
$G_b(x,y)=-u(\min(x,y);b)v(\max(x,y);b)/C_N(b)$.
The minus sign makes the derivative jump of $-\cos x\,\partial_xG$
equal to one. Differentiating \eqref{eq:u} with respect to $b$ and
pairing with $v$ gives
\[
\bigl[\cos x\,(u_bv'-u_b'v)\bigr]'=-uv.
\]
The boundary values are zero at $-a$ and $C_N'(b)$ at $a$, so
$C_N'(b)=-\int_Iuv\,\dd x$.

For real $b>\alpha$, the resolvent is positive and trace class.
On closed interior subintervals its continuous positive kernel has
trace equal to the integral of its diagonal. Exhausting $I$ by those
subintervals and using trace convergence for a positive trace-class
operator gives $\Tr(A_N+b)^{-1}=\int_I G_b(x,x)\,\dd x$.
This proves the first equality there. Both sides extend meromorphically,
so the equality holds throughout their common domain.
The series equality is the spectral theorem with \eqref{eq:bounds}.
\end{proof}

Let $D_N(b)=\det(bI-H)$ and $R_N(b)=C_N(b)/D_N(b)$, with its
removable values filled in. The first-$N$ theorem and \eqref{eq:trace}
give the exact cancellation
\begin{equation}
\frac{R_N'(b)}{R_N(b)}
=\sum_{k=N}^{\infty}\frac1{\mu_k+b}.
\label{eq:deftrace}
\end{equation}
In particular $R_N$ is nonzero on $(-\mu_N,\mu_N)$.
There is no unspecified exponential factor: \eqref{eq:trace} fixes
the logarithmic derivative directly. Related Wronskian determinant
methods are classical \cite{lesch}; no regularized determinant is needed
in this argument.

\section{Exact positive expansion and truncation bounds}
\begin{theorem}[Spectral-tail formula]
Define, for $|b|<\mu_N$,
\begin{equation}
F_N(b)=2\sum_{k=N}^\infty\artanh\frac b{\mu_k}
=\sum_{k=N}^\infty\log\frac{\mu_k+b}{\mu_k-b}.
\label{eq:F}
\end{equation}
The series converges locally uniformly. At every common finite eigenmode,
\begin{equation}
\log\frac{\zeta}{1-\zeta}=F_N(\beta).
\label{eq:exact}
\end{equation}
Consequently the following series converges absolutely in matrix
operator norm:
\begin{align}
\boxed{\ \mathcal P_N=2\sum_{k=N}^\infty
                  \artanh(H/\mu_k)\ }
&=2\sum_{r=0}^\infty
\frac{Z_N(2r+1)}{2r+1}H^{2r+1}, \label{eq:operator}\\
Z_N(s)&=\sum_{k=N}^\infty\mu_k^{-s},\qquad s\ge1.
\nonumber
\end{align}
Every coefficient on the right is strictly positive.
\end{theorem}
\begin{proof}
We first recover the finite concentration ratio from the normalized
connections. At a selected root $b=\beta$, the Fourier solution divided
by its nonzero left endpoint equals $u(x;\beta)$ and has right endpoint
$h=\pm1$. On the complementary arc, translation by $\pi$ gives the
same equation with parameter $-\beta$. The centered Fourier function
satisfies $U(x+2\pi)=\sigma U(x)$, where $\sigma=(-1)^{N-1}$.
Thus the complementary normalized endpoint value is $\sigma h$.
Parseval and \eqref{eq:norm} imply
\[
\frac{\zeta}{1-\zeta}
=\sigma\frac{C_N'(\beta)}{C_N'(-\beta)}
=\frac{R_N(\beta)}{R_N(-\beta)}.
\]
The last equality uses $D_N(-b)=(-1)^ND_N(b)$ and its derivative;
it also covers the central zero for odd $N$ by removable values.

Integrating \eqref{eq:deftrace} from $-b$ to $b$ gives \eqref{eq:F}
as $\log|R_N(b)/R_N(-b)|$. The zero-free interval fixes the positive
ratio, and $\sum_{k\ge N}1/\mu_k<\infty$ justifies integration term by
term. This proves \eqref{eq:exact}.
Because $h_N/\mu_N<1$, the scalar artanh expansions converge uniformly
on the whole finite spectrum. Their absolute norm sums are bounded by
$2\sum_{k\ge N}\artanh(h_N/\mu_k)<\infty$.
Functional calculus and absolute convergence prove \eqref{eq:operator}.
\end{proof}

The function $F_N$ is odd and strictly increasing. For $0<b<\mu_N$,
$F_N(b)/b$ is strictly increasing, by its positive Taylor coefficients.
In particular $\operatorname{sgn}(2B-I)=\operatorname{sgn}H$,
including the zero mode. For $N\ge2$, \eqref{eq:operator} also gives
zero entries between sites of equal parity and strictly positive
entries between sites of opposite parity: every odd power of the
positive-bond matrix $H$ is entrywise nonnegative, and an odd-length
path connects every opposite-parity pair. The signs reverse for $K_N$.
These are finite-$N$ conclusions. The leading scaling-limit power
series in \cite{eisler} is a related result with different coefficients.

\begin{corollary}[Explicit omitted-spectrum bound]
Let $0\le b\le h_N$ and let $M\ge N$ be the first omitted index.
If $\rho=(\alpha+b)/[cM(M+1)]<1$, then
\begin{equation}
\frac{4b}{M}\ \le\
F_N(b)-2\sum_{k=N}^{M-1}\artanh(b/\mu_k)
\ \le\ \frac{2b}{c(1-\rho)M}.
\label{eq:tailbound}
\end{equation}
For negative $b$, oddness gives the corresponding signed bounds.
The upper bound with $b=h_N$ also bounds the matrix operator-norm
truncation error in \eqref{eq:operator}.
\end{corollary}
\begin{proof}
Use $z\le\artanh z\le z/(1-z)$ for $0\le z<1$.
The upper eigenvalue bound in \eqref{eq:bounds} gives the lower
estimate. For the upper estimate,
$\mu_k-b\ge c(1-\rho)k(k+1)$ for $k\ge M$.
Now $\sum_{k=M}^\infty1/[k(k+1)]=1/M$.
The norm statement follows from the spectral theorem and monotonicity.
\end{proof}
The bound concerns the exact differential eigenvalues. Computed
approximations require their own error enclosures; an ordinary
Galerkin calculation does not automatically make this a numerical certificate.

\begin{corollary}[Acceleration and a quantitative certificate]
For $0\le b\le h_N$, set $g_b(\mu)=2[\artanh(b/\mu)-b/\mu]$.
Then
\begin{equation}
F_N(b)=2bZ_N(1)+\sum_{k=N}^\infty g_b(\mu_k).
\label{eq:accelerate}
\end{equation}
Let $M\ge N$, $d=c-\alpha/[M(M+1)]>0$, and
$q=b/[dM(M+1)]<1$. The omitted sum obeys
\begin{equation}
0\le \sum_{k=M}^\infty g_b(\mu_k)
\le \frac{2b^3}{15(1-q^2)d^3M^5}.
\label{eq:fasttail}
\end{equation}
If $|\widehat Z-Z_N(1)|\le\delta_Z$ and
$|\widehat\mu_k-\mu_k|\le\delta_k$ for $N\le k<M$, with
$a_k=\widehat\mu_k-\delta_k>b$, the absolute error of the resulting
finite approximation is at most
\begin{equation}
2b\delta_Z+
\sum_{k=N}^{M-1}\frac{2b^3\delta_k}{a_k^2(a_k^2-b^2)}
+\frac{2b^3}{15(1-q^2)d^3M^5}.
\label{eq:cert}
\end{equation}
\end{corollary}
\begin{proof}
For $0\le z\le q<1$, the positive series gives
$\artanh z-z\le z^3/[3(1-q^2)]$.
For $k\ge M$, \eqref{eq:bounds} gives $\mu_k\ge d\,k(k+1)$.
Moreover
\[
\frac1{[k(k+1)]^3}\le\frac15
\left(\frac1{k^5}-\frac1{(k+1)^5}\right),
\]
since $(k+1)^5-k^5-5k^2(k+1)^2=5k^2+5k+1>0$.
Summing proves \eqref{eq:fasttail}.
Finally $|g_b'(\mu)|=2b^3/[\mu^2(\mu^2-b^2)]$ decreases for $\mu>b$.
The mean-value theorem and the triangle inequality give \eqref{eq:cert}.
\end{proof}

\section{Center law: exact structure and formal predictions}
Let $J=2H/N$, $Jv=\lambda v$, and $t=\lambda/N=2\beta/N^2$.
Write
\[
\varepsilon=\log\frac{1-\zeta}{\zeta},\qquad
e=\varepsilon+N\eta(t),\qquad
\eta(t)=\pi t\,{}_3F_2(1/4,1/2,3/4;1,3/2;4t^2).
\]
The exact involution sends $(t,\varepsilon,e)$ to
$(-t,-\varepsilon,-e)$. Thus a leading relation involving
$e/|\varepsilon|$ cannot have the same positive sign on both branches.
Symmetry alone does not fix the sign or magnitude of $e/\varepsilon$.

The previously reduced formal coefficients \cite{periods} give
\begin{align}
\eta(t)&=\pi t+\tfrac\pi4t^3+O(t^5),&
\varphi_1(t)&=\tfrac\pi4t+\tfrac{25\pi}{32}t^3+O(t^5),\nonumber\\
\varphi_3'(0)&=-\tfrac{17\pi}{64},&
\varphi_5'(0)&=\tfrac{169\pi}{256}.
\label{eq:formal}
\end{align}
Therefore the proposed expansion
$\varepsilon=-N\eta+\varphi_1/N+\varphi_3/N^3+\cdots$
has the sign-consistent formal center law
\begin{equation}
e\sim\frac{\pi\lambda}{4N^2}
\sim-\frac{\varepsilon}{4N^2}.
\label{eq:center}
\end{equation}
This coefficient calculation does not prove a plunge asymptotic:
a fixed-interior expansion away from $t=0$ cannot simply be evaluated
along a sequence of modes tending to zero.

The tail representation gives an exact, distinguished interpolation
with
\begin{equation}
F_N'(0)=2Z_N(1).
\label{eq:slope}
\end{equation}
Differentiating the \emph{formal candidate} at zero consequently
suggests the separate trace asymptotic
\begin{equation}
Z_N(1)\ \mathrel{\stackrel{\rm formal}{\sim}}\
\frac\pi N-\frac{\pi}{4N^3}
+\frac{17\pi}{64N^5}-\frac{169\pi}{256N^7}+\cdots .
\label{eq:traceformal}
\end{equation}
Neither the spectral period conjecture at discrete points nor a
value-only remainder estimate alone justifies differentiating an
interpolant. Equation \eqref{eq:traceformal} is an additional formal
prediction to be tested and proved independently.
For even $N$, it can be tested directly through
$Z_N(1)=C_N'(0)/C_N(0)$, since $D_N$ is even and nonzero at zero.
For odd $N$, the removable-value formula is
$Z_N(1)=C_N''(0)/(2C_N'(0))$.

\paragraph{One further exact formal coefficient.}
The linear center coefficient can be continued one order without
assuming a spectral asymptotic. Set $p=\cos x$,
$q=(p-2t)/(4p)$, $r=1/(4p^2)$, and solve the formal equation
\begin{equation}
W^2+\tfrac12(\log W)''-\tfrac14[(\log W)']^2=N^2q+r,
\quad
W=N\sqrt q\,(1+a_1N^{-2}+a_2N^{-4}+a_3N^{-6}+a_4N^{-8}+\cdots).
\label{eq:Riccati}
\end{equation}
Primes in \eqref{eq:Riccati} mean $x$ derivatives.
Coefficient comparison in the ring $\mathbb Q(p)[t]/(t^2)$, using
$\partial_x^2 f=(1-p^2)f_{pp}-pf_p$, gives
\begin{align}
a_4(p,0)&=
\frac{4096p^6-143616p^4+493824p^2-375733}{128p^8},\nonumber\\
\partial_ta_4(p,0)&=
-\frac{16p^8-50952p^6+825534p^4-2268425p^2+1553721}{16p^9}.
\label{eq:a4}
\end{align}
The recursion fixes each new coefficient linearly, so these expressions
are checked by substitution after the three preceding orders.
For the counterclockwise period convention
$\varphi_7=i\oint\sqrt q\,a_4\,\dd x$ with
$\dd x=-\dd p/\sqrt{1-p^2}$, a fixed circle $2|t|<R<1$ yields
\begin{equation}
\varphi_7'(0)
=\pi\mathop{\rm Res}_{p=0}
\frac{\partial_ta_4(p,0)-a_4(p,0)/p}{\sqrt{1-p^2}}
=-\frac{50777\pi}{16384}.
\label{eq:phi7}
\end{equation}
This is an exact formal coefficient identity. The associated
$N^{-9}$ coefficient suggested for \eqref{eq:traceformal} is
$50777\pi/16384=9.7363678083\ldots$; that spectral asymptotic is
not proved by the residue calculation.

\section{Scope and verification}
Direct finite-matrix center-mode calculations at
$N=32,64,128,256,512,1024$ give
$-4N^2e/\varepsilon$ increasing from $0.9993955762$ to
$0.9999992990$ at 100-digit working precision.
An independent connection calculation of $Z_N(1)$ at
$N=8,16,32,64,128$ gives the residual after the four displayed terms
in \eqref{eq:traceformal}, multiplied by $N^9$, as
\[
8.7298271,\quad9.4540545,\quad9.6634554,\quad
9.7179853,\quad9.7317624.
\]
These values are consistent with the coefficient in \eqref{eq:phi7}.
At $N=128$, increasing precision from 130 to 180 digits and the
Taylor degree from 600 to 800 changes the retained trace by less than
$3.88\times10^{-90}$. These are numerical checks, not error enclosures
or proofs of the limits.

The proofs above identify an exact finite reference function and
give convergent spectral representations. They do not establish
the large-$N$ expansion of the remaining differential spectrum, the
uniform quantum-period remainder, the full Wilson relative-entropy
law, or a field-theory mass gap.
The decomposition isolates a further precise target: estimate the
positive traces $Z_N(2r+1)$ with uniform control sufficient to
compare \eqref{eq:operator} to the declared period candidate.
The reciprocal-trace, positive-moment and whole-line results supply the
stated fixed-interior estimates. For the specified infinite-chain Wilson
comparison at fixed positive scaled mass, Theorem 1 of \emph{Finite Part
of Wilson--Dirac Relative Entropy} gives the finite part and an
\(O_\mu(1/\log L)\) remainder; Theorem 3 of \emph{Uniform Modular
Control and Nonlinear Convergence for Wilson--Dirac Intervals} supplies
the even-dimensional bound \(\|E_N\|\le7\). A moving-edge expansion
remains open. Neither the finite-part theorem nor that norm bound
promotes the formal coefficient calculation above to an additional
spectral asymptotic.
Separate reproduction material records numerical discretizations,
precision comparisons, formal coefficient checks and their limitations.

\paragraph{Acknowledgments and provenance.}
This work develops the finite reference problem in
\emph{From Prephysical Hypotheses to Physical Models}.
The earlier half-filled connection proposal was supplied in
correspondence prepared with Claude (Anthropic).
The present spectral-tail argument and checks were developed with
ChatGPT (OpenAI). The wider program has also used Gemini (Google),
Grok (SpaceXAI), and DeepSeek (High-Flyer); no particular proof here
is attributed to those systems. Assistance implies neither
external peer review nor endorsement. No priority claim is made.

\begin{thebibliography}{9}
\bibitem{eisler} V. Eisler and I. Peschel, Analytical results for the
entanglement Hamiltonian of a free-fermion chain,
\emph{J. Phys. A} \textbf{50} (2017), 284003,
\href{https://arxiv.org/abs/1703.08126}{arXiv:1703.08126},
especially equations (3), (6)--(10), (21)--(22).
\bibitem{kaper} H. G. Kaper, M. K. Kwong and A. Zettl,
Characterizations of the Friedrichs extensions of singular
Sturm--Liouville expressions,
\emph{SIAM J. Math. Anal.} \textbf{17} (1986), 772--777.
\bibitem{littlejohn} L. L. Littlejohn and A. Zettl,
The Legendre equation and its self-adjoint operators,
\emph{Electron. J. Differential Equations} (2011), No. 69, 1--33,
especially Theorem 3.6.
\bibitem{lesch} M. Lesch, Determinants of regular singular
Sturm--Liouville operators,
\href{https://arxiv.org/abs/math/9902114}{arXiv:math/9902114}.
\bibitem{periods} K. Justus,
\emph{Higher-Order Modular Periods for the Half-Filled Sine Kernel},
research manuscript, October 2026.
\end{thebibliography}
\end{document}
